\section{Introduction}

Rocket engines require high pressures to operate efficiently. Small scale liquid propulsion such as satellite propulsion and amateur liquid rocket propulsion runs almost entirely on pressure fed feed systems due to their simplicity and reliability. Larger orbital rockets use turbopumps because the mass of the required tanks becomes prohibitive. However, as space launch becomes exponentially more accessible, there is a growing interest in small scale rocket propulsion systems. For example, lunar landers, satellite propulsion and amateur liquid rockets are seeing a resurgence in interest. Due to the historical focus on large turbopumps for such applications, turbopumps at the smaller scale are not as well documented and have different design considerations. A turbopump is a pump and the turbine that drives it, and both are punished by the small scale.

\subsection{Small scale rocket pump challenges}

Pump driven propulsion is the solution orbital liquid rockets use to solve the tank mass problem. By storing propellants at lower pressures and pumping them to the required pressures allows for a large reduction in tank mass. The pump is usually driven by turbines powered by burning a portion of the propellants, or by the heat generated by the cooling of the engine. In recent years, battery and electric pump technology has advanced to such a point that for smaller rockets such as the Electron rocket by Rocket Lab \cite{electron}, electric pumps are starting to become a more viable alternative to turbine driven pumps. \todo{this paragraph might be redundant?}

Pumps are generally characterized by their flow rate $Q$, the head $H$ and rotor speed $n$ \cite{gulich}. These three parameters largely decide the pump geometry and type of pump. These three parameters can be related by a quantity known as specific speed, $n_q$, $N_s$ or $\omega_s$ depending on the units and definitions used:

\begin{equation}
    n_q = \frac{n \sqrt{Q}}{H^{3/4}} \quad [\text{RPM}, \text{m}^3/\text{s}, \text{m}] \label{eq:specific_speed}
\end{equation}

\begin{wrapfigure}{r}{0.6\textwidth}
    \centering
    \includegraphics[width=0.58\textwidth]{pump_map_balje.pdf}
    \caption{Optimum pump geometry and attainable efficiency against specific speed. Adapted from Balje \cite{baljeBaljeTurbomachinesGuide1981} with specific speed in $n_q$ units. } 
    
    \label{fig:pump_type_by_specific_speed}
\end{wrapfigure}

Different pump geometries are optimal for different regions of specific speed as can be shown by \cref{fig:pump_type_by_specific_speed}. The radial centrifugal impeller reaches its best efficiency roughly between an $n_q$ of 10 and 40 \cite{gulichCentrifugalPumps2020,huzel}, and rocket engine pumps are normally sized to sit in this band, splitting the head across multiple stages when a single stage would otherwise fall below it. This works great for large rocket engines, but scaling this down to smaller engines becomes increasingly difficult. All rocket engines require high pressures, but smaller rocket engines require less flow rate. To stay in the specific speed according to \cref{eq:specific_speed} that allows for high efficiency in the centrifugal pump requires increasing rotational speed.


\todo{Caption should say adapted from sp 8107 to use nq units, design point also not yet marked}

To achieve the specific speed $n_q$ of \valSamplepumpNq for a smaller rocket engine of around \valSamplepumpThrust requiring \valSamplepumpMdot of kerosene at \valSamplepumpP, \cref{eq:specific_speed} leads to rotational speeds of around \valSamplepumpN. Such high rotational speeds leads to challenging mechanical design conditions. To solve this problem, a single pump can be separated into multiple stages, separating the head rise to increase the specific speed of each stage. This is a common solution to pumping hydrogen. This thesis however addresses the other solution, to study a single stage pump geometry at this low specific speed. In rocket applications, the lower component efficiency could trade better by other factors such as lower mass, cost, reliability and simplicity.

As shown in \cref{fig:pump_type_by_specific_speed}, at lower specific speeds, the partial emission pump, of which the Barske impeller is the main form, becomes the optimum efficiency pump geometry at specific speeds of 5-10. The term partial emission refers how the pump outlet only taps a small arc of the pump periphery. Full emission pumps are designed to have the impeller discharge around the whole annulus, collecting flow through a continuous volute. Literature uses the names partial emission and Barske interchangeably.

%The NASA monograph shows "Barske" and "Partial emission" as separate lines due to older classification   
However, even though these pumps are the most efficient at these low specific speeds, they physically cannot achieve the high efficiencies of a full emission pump. This is the main limitation of trying to operate at low specific speeds. Due to its lower efficiency and the lack of pump fed liquid rockets for small engines, there is a lack of literature, design knowledge and experimental data surrounding these pumps. This thesis aims to address this gap by reviewing existing literature, designing and manufacturing a partial emission centrifugal pump and comparing its experimental performance specifically in the use case of rocket engines.

The lack of data is specifically a lack of experimental data. Open studies of this regime do exist, but they are mostly numerical. A Barske type impeller at $n_q = 4.7$ was studied with CFD and reached an efficiency of 31.5\% \cite{knyazevaDesignFluidDynamics2021}. A conventional closed impeller was designed at the same specific speed as this thesis, but it was CFD only, unvalidated, and required closed channels of 0.31 mm width \cite{kimDevelopmentUltraLowSpecific2022}. A partial emission rocket e-pump for a lunar lander was designed for 630 m at 40,000 RPM, again in CFD \cite{danieliConceptualDesignElectric2023}. What is scarce is open experimental characterisation at this specific speed, with a measured loss breakdown and measured cavitation behaviour, and that is what this work supplies. A NASA trade study of low thrust rocket pumps covers the same trade and carries the Barske pump as a candidate \cite{macgregorRotatingPositivedisplacementPumps1974}.

%Balje's pump selection diagrams show why this corner of the design space is poorly covered, as the efficiency islands collapse below a non-dimensional specific speed of around 0.1, right where this duty point sits \cite{baljeBaljeTurbomachinesGuide1981}. this doesnt seem to add anything, page 6 of the figure here https://barber-nichols.com/wp-content/uploads/2019/11/how_to_select_turbomachinery_for_your_application.pdf

\subsection{Small scale rocket turbine challenges}

On the other end of the turbopump shaft is the turbine which drives the pump. The turbine faces the same problem as the pump. If using the gas generator cycle, the small turbine leads to lower blade speeds $u$. To extract the maximum work from the gas, gas velocity $c_s$ must be high, leading to a supersonic blade turbine. The high gas velocity relative to blade speed leads to the turbine running at a lower efficiency than the optimum $u/c_s \approx 0.5$ \cite{findsource}. The smaller mass flow also cannot fill the annulus at a sensible blade height, leading to partial admission where only a fraction of the annulus has gas flowing through it at any given time. This makes the stage lossy and awkward to design. In addition to these design limitations, supersonic turbine design is far less studied than subsonic, where methods are scattered across older sources. \todo{need to address subsonic turbines?}

This thesis tackles the design of a turbopump in the context of small scale liquid rocket engines which forces both pump and turbine off optimum design conditions, then demonstrates their coupled performance in a full integration test. The thesis is framed around the design of a \valPumpRpmDesign turbopump using water as an inert simulant for the pumped rocket propellants. Then, a supersonic partial admission impulse turbine is designed to drive the pump using compressed air as the working fluid. The two are then coupled together and tested as a single system to evaluate performance.

\subsection{Scope and Limitations}

This thesis covers the theory, design, manufacture and testing of a small scale turbopump system on inert propellants. The pump and turbine were tested up to \valPumpRpmMaxreached, significantly lower than the design speed of \valPumpRpmDesign in the first and final iterations of the conducted experiments. This means the data validation can only happen for the off-design point and extrapolation to the design point is required. Cryogenic and hot gas design and testing are strictly outside the scope of this thesis. Operational safety procedures are included in the appendix.

\subsection{Aims and Objectives}
\label{sec:intro_objectives}
The aim of this thesis is to investigate the performance of a small scale rocket low specific speed turbopump with a Barske pump and a supersonic partial admission turbine. To this end, the objectives of this thesis are:

\begin{enumerate}
    \item Pump: Assemble a Barske pump design methodology from the literature, use it to design and build the pump, then characterise it experimentally for head-flow, efficiency and cavitation, and compare the measurements against the methodology's predictions and prior literature.
    \item Turbine: Assemble a design methodology for a single-stage partial-admission supersonic impulse turbine from the literature, use it to design and build the rotor, then validate those models against the achievable test data.
    \item Turbopump: Couple the pump and turbine and demonstrate them running as a single machine, then use that result to check whether the combined methodology predicts a working turbopump of this class.
\end{enumerate}

\subsection{Thesis Structure}

% old nine-chapter paragraph commented out (restructured to six chapters, Knoll 06-12):
% Chapter 2 establishes the pump theoretical background: pump specific speed, Barske impeller characteristics, impulse turbine velocity triangle theory, and the relevant prior literature. Chapter 3 details the pump design and sizing methodology. Chapter 4 covers turbine design, including structural analysis of the rotor at \valPumpRpmDesign. Chapter 5 describes the custom instrumentation and data acquisition system. Chapter 6 documents the test rig and assembly. Chapter 7 presents the experimental results. Chapter 8 discusses the findings in the context of design predictions and prior literature. Chapter 9 draws conclusions and recommends future work.
Chapter 2 establishes the theoretical background for both machines: the Barske pump and its sizing and loss models, and the supersonic impulse turbine theory including blade profile design and loss models. Chapter 3 applies these models to design the actual pump and turbine, including the mechanical and structural design. Chapter 4 describes the experimental setup: the test configurations, the water and air circuits, and the custom instrumentation and data acquisition system. Chapter 5 presents the experimental results and discusses them against the design predictions and prior literature. Chapter 6 draws conclusions against the objectives and recommends future work.